% Chapter 4: Data, labels and the dataset audit. Draft 23 Nov 2026. Every number is traced in ch04_data_NOTES.md.
% Tables are in the appendix (thesis_tables.tex); they are referenced only by \ref{tab:...}.
\chapter{Data, labels and the dataset audit}\label{ch:data}

This chapter describes the images, the cell sets, the label sets and the splits used in the thesis, and then the audit of the
development pool that we ran before fixing the protocol. The audit comes before the protocol on purpose: its findings on the
line-width convention, on split integrity and on the label release are the reasons for the rules of Chapter~\ref{ch:protocol}.
All tables of this chapter are in the appendix: the cell sets in Table~\ref{tab:T1}, split integrity in Table~\ref{tab:T2}, the
label sets in Tables~\ref{tab:T3a}, \ref{tab:T3b} and~\ref{tab:T3c}, and the audit in Tables~\ref{tab:T19a} and~\ref{tab:T19b}.
Table~\ref{tab:T3b} counts label changes; it is unrelated to the few-shot track of Chapter~\ref{ch:t3b}.

Several model recipes appear here before Chapter~\ref{ch:methods} describes them. PF is the baseline: U-Net++
\citep{zhou2018unetpp} with an ImageNet ResNet-34 encoder and an encoder-preserving fine-tune (Chapter~\ref{ch:legacy}); HOUSE is
the same model with a plain fine-tune; HRNET replaces the encoder by HRNet-W18 \citep{wang2021hrnet}; SWADBN is PF with weight
averaging; C3X combines the last two. The audit scores the out-of-fold predictions of these five recipes at five seeds each, 25
models per pool cell.

\section{Images, sources and cell sets}\label{sec:data-sets}

Every image is a grayscale electroluminescence (EL) image of 512 $\times$ 512 px that shows one cell and parts of its neighbours
(Section~\ref{sec:data-other}). The main labelled set is BenchmarkEL \texttt{dataset\_20211104} \citep{pratt2023benchmark}: 582 cells
in a legacy split of 478 train, 54 validation and 50 test cells, with masks over 24 classes (12 structural classes and 12 defects).
The cells come from five sources with different image contrast and label conventions. In the development pool the counts are
ARTS 182, BMRK 50, CFVS 169, CSIR 114 and SDLE 17 cells. BMRK cells are cells of the ELPV set \citep{deitsch2019elpv}, downscaled and
padded. Over the three BenchmarkEL releases and the E-SCDD extension \citep{hanifi2025escdd}, the labelled corpus holds 816 real
labelled cells, 787 unique physical cells and 463 physical-module groups; the 93 synthetic E-SCDD files are excluded everywhere.

Table~\ref{tab:T1} lists the sets built from this corpus; Figure~\ref{fig:F4} shows how they relate.
\begin{itemize}
  \item \textbf{Pool P532} (legacy train and validation, 478 + 54 cells). All development, the five-fold cross-validation (CV5) and
    the full-pool training of the final models use these cells only.
  \item \textbf{test50} (the legacy test split; 50 cells from 48 image ids). The cells were hand-picked: 48 hold crack under v1 (46
    under v1.1, Section~\ref{sec:data-labels}) and 16 hold inactive area, and 3 cells hold 59\% of the inactive pixels. The set is sealed. It was scored in every experiment before 24 Oct
    2026 without being counted, and opened 19 counted times from 24 Oct to 21 Nov 2026, the last time for a public table that
    decides nothing. Section~\ref{sec:data-split} shows that test50 overlaps the pool, so every test50 number in this thesis is an
    in-distribution number with known leakage and serves as context only.
  \item \textbf{E1} (106 cells of release 20221008 that are neither among the 582 cells nor in the unlabelled pool). The cells come
    from 71 image ids, 66 CSIR and 5 BMRK; 101 of the 106 cells are CSIR. 62 cells are half-cut cells, a layout that the pool does not
    contain at all; 44 are full cells. Crack is labelled in 32 cells and inactive area in 10. E1 is our test of unseen modules
    (supported, not proven; Section~\ref{sec:data-split}). It decided the final evaluation (Chapter~\ref{ch:e1_result}); afterwards it
    served as training and test data of the few-shot track (Chapter~\ref{ch:t3b}), so no zero-shot claim on E1 is made after that
    use. Its strata in cells and clusters are full 44 / 24, half-cut 62 / 47, crack 32 / 21 and crack-free 74 / 55.
  \item \textbf{Forbidden stems}. No training list may contain any of 167 stems: the 50 legacy test cells and the 117 cells first
    listed for E1 from release 20221008 (E1 is the 106 of them that are not in the unlabelled pool).
  \item \textbf{Other sets.} Of the 117 cells newly annotated in E-SCDD, 51 full cells and 36 half cells served as development
    validation sets only.
    The unlabelled pool of 11,938 CSIR and ARTS crops served the semi-supervised work of Chapter~\ref{ch:legacy}.
\end{itemize}

Sealing means that no development step reads a label or a prediction of a test50 or E1 cell; the duplicate search of
Section~\ref{sec:data-split} reads sealed images as pixels only and writes pair lists. On terms of use: BenchmarkEL has no licence
file, and the BMRK cells are ELPV images under CC BY-NC-SA 4.0, so every model trained on the pool is research-only as trained
(Chapter~\ref{ch:protocol}).

\FIG{F4}{Cell sets and splits (counts as in Table~\ref{tab:T1}). The pool P532 is split twice: cv5v2 (2-seed screens, 330 groups,
with a module leak of 24 modules, 91 cells) and cv5v3 (5-seed rows, 285 groups, none spans folds; the 416 clean cells are read beside
all 532); bars give held-out fold sizes. Sealed: E1 (106 cells, 71 image-id clusters, 101 CSIR, 62 half-cut) and test50 (50 cells;
in-distribution, known leakage, context only). The 167 forbidden stems may enter no training list.}

\section{Label sets}\label{sec:data-labels}

Four label sets are in use (Table~\ref{tab:T3a}).
\begin{itemize}
  \item \textbf{v0} is the 20211104 release as published, values 0 to 23 and no ignore value. It is the label set of the published
    test50 numbers \citep{pratt2023benchmark}; we use it only for the test50 public table.
  \item \textbf{v1} is the E-SCDD mask (577 cells) with classes $\geq$ 24 set back to the 20211104 label, plus the 20211104 mask for the 5
    cells that E-SCDD lacks. The pre-protocol work (Chapter~\ref{ch:legacy}) trained and scored on v1 unless stated.
  \item \textbf{v1.1} is v1 with every scuff or star pixel set to ignore (255) and one BMRK cell given back its 20220723 mask; in v1 that
    cell carried the mask of another cell. Ignore pixels are not scored. v1.1 is the training and scoring label set of every CV5 read
    and of both full-pool models.
  \item \textbf{The E1 set} is the release 20221008 masks with class 24 (sp mono of the newer releases) mapped to 2 and classes 25 to 28
    (scuff, corrosion cell, brightening, star) mapped to ignore. E1 labels are about 1.5 $\times$ wider than the pool labels.
\end{itemize}

The change from v1 to v1.1 is small (Table~\ref{tab:T3b}): 16 of 478 train cells change, with 3,484 px set to ignore (3,246 scuff,
238 star); no validation cell changes; 2 test cells change, with 236 px. Train crack cells go from 209 to 206; inactive labels are
untouched. The reason for the change was that v1 labels scuff streaks as crack on 649 px in 6 train cells, and scuff streaks are a
confounder on the half-cut cells of E1. A cell counts as crack-free under the scoring label set; on v1.1 that is 303 of 532
pool cells and 226 of 416 clean cells.

The label set matters at scoring time, not at training time (Table~\ref{tab:T3c}). Training PF on v1 instead of v1.1 (two seeds,
both scored on v1.1) gives relF1@2 $-0.0014$ $[-0.0099, +0.0064]$, strict crack IoU $+0.0004$ $[-0.0058, +0.0069]$, inactive IoU
$-0.0140$ $[-0.0690, +0.0362]$ and mIoU-R13 $-0.0020$ $[-0.0081, +0.0038]$: no pooled metric changes by a measurable amount at two
seeds. Scoring is a different matter. Between the 20211104 release and v1.1, 169 of the 532 pool cells change and the crack pixels
grow from 260,084 to 307,546 (+18\%). The same 25 out-of-fold models score PF at strict crack IoU 0.320 against 0.365 (+0.045) and
relF1@2 0.822 against 0.846 (+0.025); inactive is unchanged, and contrasts between arms move by at most 0.008, with no pass changed.
The models were trained on v1.1, which favours v1.1. For scale, the C2 recipe of this thesis gains $+0.0222$ relF1@2 over PF on CV5
(selection-biased; Table~\ref{tab:T7a}), less than the release shift. Our reading is that a published crack level means little
unless its release and label set are stated.

\section{Module identity and split integrity}\label{sec:data-split}

Cells of one physical module share grain patterns, and neighbouring crops of one module image share area, so a cell can be
predicted well because its module was trained on. The first five-fold split, cv5v2, grouped the pool into 330 module groups. An audit of physical-module identity over all 700 image ids on disk (pool,
sealed and unlabelled) tested this grouping. It used two kinds of evidence. On ARTS images, the printed serial numbers were compared
as text strips. For all sources, a masked local-feature matcher (SIFT keypoints with a RANSAC similarity fit) linked two image ids when a cell pair
matched under a whole-cell grid shift with enough inliers. The thresholds were written down before the full run and set on a
calibration subset, so their error rates are in-sample. In the full run, 20 of 20 same-serial pairs were also matched by pixels
and 0 of 2,906 different-serial pairs were. The 700 image ids form 561 physical modules under the rule we use (527
under the rule fixed in advance, which over-links some image ids that hold only unlabelled crops). Our rule was chosen after the
full run, but it changes nothing for pool or sealed cells. For BMRK and SDLE, which have one image id per cell, module identity
cannot be decided.

\paragraph{The cv5v2 leak.} 24 physical modules span two cv5v2 folds. Their 91 pool cells (ARTS 48, CFVS 32, CSIR 11; 22.7\% of
the crack pixels) are held out while a cell of the same module trains: for 4 of them the same physical cell, for 38 a neighbouring
crop, for 49 the same module without shared area (Table~\ref{tab:T2}). These cells score higher than comparable cells: within the
same source and cell type, the affected cells exceed the unaffected ones by $+0.072$ crack IoU and $+0.059$ relF1@2. The
comparison is descriptive, not causal. The affected modules were imaged several times and many carry large stress-test cracks, and
the gap does not grow with closeness: cells whose neighbouring crop trains score lowest, which a pure leak would not produce. If
the whole gap were leakage, the 532-cell headline would be too high by about 0.016 crack IoU and 0.013 relF1@2. Scoring without
the 91 cells lowers PF from 0.365 to 0.347 strict crack IoU and from 0.844 to 0.826 relF1@2, and no contrast changes sign.

\paragraph{The corrected split cv5v3.} We regrouped the pool before any five-seed run: cv5v2 groups, the physical modules of the
audit, four open image-id pairs, earlier links and the SDLE modules named in the file stems give 285 groups, none of which spans
folds. 47 cells move; the held-out folds hold 107, 110, 103, 101 and 111 cells. BMRK cells stay grouped by image id, a stated limit,
and SDLE cells are held out in folds 0, 2 and 4 only. Every cv5v3 read is printed on all 532 cells and on the 416 \emph{clean}
cells: 116 cells are excluded (the 91 affected cells, the 8 cells of the open pairs and the 17 SDLE cells), and the 416 keep their
fold in both splits.

\paragraph{The leakage read.} A difference-in-differences read, registered before any cv5v3 run, compared the change of the pooled
score from cv5v2 to cv5v3 on the 69 excluded (non-clean) cells that kept their fold with the same change on the 416 clean cells. For PF at five seeds it
gives $-0.0041$ $[-0.0136, +0.0069]$ strict crack IoU and $+0.0026$ $[-0.0092, +0.0174]$ relF1@2. Every interval of both arms (PF
and HOUSE) and both strata holds 0, so by the registered rule the leak is \emph{not separated}. A leak of the size that the module audit
allowed as an upper bound ($+0.07$ crack IoU on the affected cells) lies outside the PF intervals; by the registered rule this is
a description of the intervals, not a conclusion. Inside the pool, 61 images also overlap a pool image of the adjacent
position of the same module (35 pairs, 16\% of the crack pixels), so parts of one crack are scored twice; both images sit in one
cv5v3 group, so this is no fold leak.

\paragraph{The sealed sets.} No E1 image id is linked to any other image id, and the largest inlier count between an E1 cell and a
pool cell is 7, which is noise level. The test covers about a third of the image-id pairs: the half-cut E1 modules have no layout
counterpart in the pool, and BMRK identity cannot be tested. ``Unseen modules'' is therefore supported by the pixels on disk, not
proven. test50 is in a different position. 8 test cells are the same physical cell as a pool cell (7 confirmed by pixels, 1 resting
on the earlier cell-table link only, neither confirmed nor refuted), 7 more are neighbouring crops of pool cells, and 19 share a
module with pool cells without shared area; only 16 test cells lie in image ids without any link. test50 numbers are
in-distribution numbers with known leakage and are used as context only.

\section{Strict crack IoU measures the line-width convention}\label{sec:data-f1}

Crack labels are lines of 1 to 2 px drawn on a dark line about 4 px wide (Chapter~\ref{ch:background}). The audit measured what
this does to strict crack IoU (fault F1 in Table~\ref{tab:T19a}). Against itself, a crack label scores IoU 0.483 when shifted by
1 px, 0.249 when shifted by 2 px, 0.465 when reduced to its centreline and 0.605 when redrawn 3 px wide, pooled over the 229 crack
cells; relF1@2 is 1.0 in every case. We then re-scored the same stored predictions of the 25 models against three versions of the
same crack labels: the 1-px centreline of v1.1, v1.1 itself, and the centreline drawn 3 px wide. PF scores strict crack IoU 0.233,
0.365 and 0.414, while relF1@2 stays at 0.824, 0.846 and 0.849.

The convention also moves contrasts between arms (Table~\ref{tab:T19b}, Figure~\ref{fig:F5}). HOUSE minus PF in strict crack IoU is $+0.0179$ and
passes on v1.1, falls to $+0.0015$ without a pass on centreline labels and rises to $+0.0367$ on 3-px labels. SWADBN minus PF behaves
the same way: $+0.0180$ (pass), $+0.0047$ (no pass), $+0.0342$. HRNET minus PF changes sign in strict IoU, from $+0.0008$ on v1.1 to
$-0.0109$ on 3-px labels. relF1@2 contrasts move by at most 0.0031 under every variant, but one of them changes its verdict: HRNET
minus PF falls from $+0.0104$ (pass) to $+0.0073$ (no pass) on 3-px labels. C3X minus PF passes under all four variants. The
\emph{convention band} of $\pm$0.019, the largest shift of a strict-IoU contrast under the three conventions, comes from this
re-scoring; the protocol prints strict crack IoU only with this band and lets it decide nothing (Table~\ref{tab:T4}). The variants
are synthetic conventions, not a second annotator, and the 3-px variant applies BMRK's convention everywhere. We read the
strict-IoU gains of the arms that draw wider cracks (HOUSE and the weight-averaged arms) as convention effects, and relF1@2 as the
crack metric that should decide.

\FIG{F5}{Strict crack IoU and relF1@2 under four crack label conventions: the same predictions (25 out-of-fold cv5v3 models, 532 pool
cells, argmax) scored on the 1~px centreline, the 20211104 release, v1.1 and labels redrawn 3~px wide. (a) PF levels. (b, c) HOUSE,
HRNET, SWADBN and COMBO (C3X) against PF in strict crack IoU and in relF1@2, with 95\% two-stage group-bootstrap intervals; dashed
line M = 0.010. Across conventions the strict-IoU contrasts shift by up to 0.019, the relF1@2 contrasts by at most 0.0031.}

\section{Other faults of the pool}\label{sec:data-other}

\paragraph{The images are windows (F3).} ARTS, CFVS and CSIR images show one centre cell, 25 to 30\% of the non-padding pixels, plus
parts of the eight neighbouring cells, and the labels cover the whole window; BMRK and SDLE images are single cells with 75\%
padding. 57\% of the v1.1 crack pixels and 71\% of the inactive pixels lie outside the centre cell (Figure~\ref{fig:F3}). 117 of the 229 ``crack cells''
have no crack pixel in their centre cell, and 28 of the 48 ``inactive cells'' have no inactive pixel there. Defects on the
neighbour parts are harder: the 25 models recall 0.794 of crack pixels within 2 px outside the centre cell against 0.844 inside.
The centre-cell geometry was defined after the first numbers and checked by eye on 7 images. We keep the image window as the
scoring unit, so ``cell'' in every cell-level quantity of this thesis (crack-free cells, inactive area, severity) means the window.

\FIG{F3}{Pool images are windows of about 3 $\times$ 3 cells. (a to c) One pool cell each from ARTS, CFVS and CSIR, chosen by a fixed
rule, with the v1.1 crack label and the centre cell found by the audit's rule (dashed box). (d) Share of v1.1 crack and inactive pixels
outside the centre cell, per source and over all 532 pool cells (ALL includes the single-cell BMRK and SDLE images, whose shares are 0).}

\paragraph{One double annotation (F2).} The release contains one image annotated twice (a CFVS cell, identical pixels, both copies in
the pool). Mask against mask, strict crack IoU is 0.259 and relF1@2 0.829, below PF against the label (0.365 and 0.846) and the
25-model vote (0.398 and 0.868). Two images of the same physical cell without crack growth give strict IoU 0.16 to 0.37 and a
tolerant F1 at 2 px of 0.77 to 0.95. Most residual error is shared by the models: 63\% of each model's missed crack pixels are missed within 2 px
by at least 20 of the 25 models, and only 7\% by 5 or fewer. With one double annotation from one source, we treat this as a
suggestion that crack labels limit what can be measured, not as an estimate of a ceiling; shared misses also include blind spots
common to all models.

\paragraph{Inactive labels (F4).} Spurious-inactive candidates, labelled regions that almost no model predicts as inactive, cover
63 regions in 18 cells, 31,039 px or 10.9\% of the inactive pixels; ignoring them would raise PF inactive IoU from 0.504 to 0.574.
Most of these candidates are not label noise: in a blind check by AI raters (below), 30 of 40 were judged model errors and 5 of 40 label
errors, while 5 of 6 missed-inactive candidates were judged label errors. Inactive pixels sit in 33 of the 285 groups, an effective
number (Kish) of 12.9 groups, and every inactive contrast so far has an interval of about $\pm$0.05 to $\pm$0.06 (audit estimate). Conventions also
differ: 26 of 47 inactive cells have cracks drawn inside the inactive area. We report inactive results on the labels as they are,
with region counts and five seeds.

\paragraph{Crack label errors (F5).} Consensus candidates for missed cracks (found by at least 20 of the 25 models, more than 2 px
from any label) cover 128 segments in 42 cells, 5,436 px or 1.8\% of the crack pixels; candidates for spurious crack labels cover
53 components in 28 cells, 7,567 px or 2.5\%. A pre-registered check had three AI raters (three runs of one multimodal model,
calibrated on 60 planted items, each with sensitivity 29/30 and specificity 30/30) rate the 78 image-supported crack candidates
blind (examples in Figure~\ref{fig:F7}). 18 of 78 were called label errors; corrected for rater error the share is 0.244 $[0.105, 0.436]$, and most of the rest were
ambiguous at this resolution. Label errors exist, and they are a minority of the image-supported candidates. Ignoring every
candidate pixel raises PF relF1@2 from 0.846 to 0.864, but no contrast moves by more than 0.003. The detector finds only errors the
models do not share, and its recall is unknown. The raters are runs of one model, so the check measures consistency, not
independent human reliability. Related approaches estimate label errors from model confidence \citep{northcutt2021confident} or detect them per
connected component in segmentation \citep{rottmann2023labelerrors}.

\FIG{F7}{Six label-error candidates on pool cells, crack on the left and inactive on the right, one per class and majority verdict of
the three AI raters (rows: label error, model error, ambiguous; the largest candidate of each group). Each shows the image window, the
v1.1 label and the out-of-fold vote of at least 13 of the 25 cv5v3 models; the dashed box is the candidate region. The raters are runs
of one model: consistency, not independent reliability.}

\paragraph{Source conventions (F7, F8).} Pooled crack label width, in pixels per centreline pixel, is ARTS 1.59, CFVS 2.06, CSIR 2.25 and
BMRK 3.09; the source explains 29.6\% of the variance of per-cell width (Figure~\ref{fig:F6}). In every source the label is about half as wide as the
visible dark line. 10 of the 24 classes occur in one source only (splice only in ARTS, corrosion only in SDLE), so per-class IoU of
these classes measures recognition of the source. The source can also be recognised from global image statistics: a simple source
classifier reaches balanced accuracy 0.94 (chance 0.20), and 0.85 from an 8-bin power spectrum alone. This shows that a shortcut is
available, not that a model uses it. Together these findings motivate the per-source width guard of the protocol.

\FIG{F6}{Crack label width per source on the pool (v1.1), in label pixels per centreline pixel: one value per cell for the 207 cells
with at least 50 centreline pixels. Boxes show the quartiles, whiskers reach 1.5 $\times$ IQR, circles are cells beyond them; SDLE (2
cells) is drawn as points. Source explains 29.6\% of the variance of per-cell width.}

\paragraph{Few groups carry the crack contrasts (F9).} Crack pixels lie in 125 of the 285 groups, an effective number of 24.1
groups. Removing whole groups greedily, HOUSE minus PF (relF1@2 $+0.0059$) reaches 0 after 4 groups; HRNET minus PF ($+0.0104$) falls
below the margin after 1 group and reaches 0 after 17; SWADBN minus PF reaches 0 after 24 and C3X minus PF after 31
(Table~\ref{tab:T19b}). Greedy removal is a worst case, not a sampling distribution; the bootstrap interval is the latter.

\paragraph{Image faults and duplicates (F10, F11).} 2 pool images have a noise sigma above 10, 109 of 532
show an 8 $\times$ 8 block grid from earlier JPEG compression, and 19 CSIR images are under-exposed. A perceptual-hash search found
41 near-identical image pairs across groups, all ARTS or CSIR look-alikes at the same grid position, and no new leak; checks for
near-duplicates between training and test data are standard practice in benchmark audits \citep{barz2020cifair}.

\section{What the audit means for the benchmark}\label{sec:data-consequences}

Absolute levels on this dataset move by 0.02 to 0.18 with the label release, the line-width convention and the label errors, while
the relF1@2 contrasts between recipes move by at most about 0.003 under the same changes. We draw five consequences, which
Chapter~\ref{ch:protocol} makes into rules.
\begin{enumerate}
  \item No absolute crack or inactive level from this dataset should be compared with a level from another dataset or release; a
    level states its release, label set and split.
  \item relF1@2 decides crack contrasts; strict crack IoU is printed with its convention band and decides nothing.
  \item Development uses the module-grouped split cv5v3 and prints the 416 clean cells beside the 532; E1 is the unseen-module test
    (supported, not proven), and test50 is in-distribution with known leakage, context only.
  \item Width is guarded per source, because label width is a source convention; per-class IoU of one-source classes is read as
    source recognition.
  \item A cell-level quantity describes the image window. Inactive claims need region counts and five seeds, and a crack gain
    reports which module groups carry it.
\end{enumerate}
What the audit could not settle needs new data or labour: a second, independent annotation of a few dozen pool cells would turn
the double-annotation finding into a measured ceiling, and planting known errors would measure the recall of the label-error
detector.
